\documentclass[11pt,a4paper]{article}

\usepackage[utf8]{inputenc}
\usepackage[T1]{fontenc}
\usepackage[margin=2.2cm]{geometry}
\usepackage{parskip}
\usepackage{titlesec}
\usepackage{enumitem}
\usepackage{booktabs}
\usepackage{xcolor}
\usepackage{tcolorbox}
\usepackage{hyperref}
\usepackage{fancyhdr}
\usepackage{listings}
\usepackage{longtable}

\tcbuselibrary{skins,breakable}

\definecolor{accent}{RGB}{30,80,150}
\definecolor{softgray}{RGB}{245,246,248}
\definecolor{promptbg}{RGB}{250,250,252}
\definecolor{promptborder}{RGB}{200,205,215}
\definecolor{codebg}{RGB}{248,248,250}

\titleformat{\section}{\Large\bfseries\color{accent}}{\thesection}{0.8em}{}
\titleformat{\subsection}{\large\bfseries}{\thesubsection}{0.6em}{}
\titleformat{\subsubsection}{\normalsize\bfseries\itshape}{\thesubsubsection}{0.5em}{}

\hypersetup{
  colorlinks=true,
  linkcolor=accent,
  urlcolor=accent,
  pdftitle={SkillOpt Reference --- AgentIA},
  pdfauthor={AgentIA Engineering}
}

\pagestyle{fancy}
\fancyhf{}
\fancyhead[L]{\small\color{accent}SkillOpt Reference $\cdot$ AgentIA}
\fancyhead[R]{\small\thepage}
\renewcommand{\headrulewidth}{0.4pt}

\newtcolorbox{keybox}[1][]{
  colback=softgray, colframe=accent, boxrule=0.4pt,
  arc=2pt, left=8pt, right=8pt, top=6pt, bottom=6pt,
  breakable, #1
}

\newtcolorbox{promptbox}[1][]{
  colback=promptbg, colframe=promptborder, boxrule=0.4pt,
  arc=2pt, left=8pt, right=8pt, top=6pt, bottom=6pt,
  fontupper=\small\ttfamily, breakable, #1
}

\lstset{
  basicstyle=\small\ttfamily,
  backgroundcolor=\color{codebg},
  frame=single,
  rulecolor=\color{promptborder},
  breaklines=true,
  columns=fullflexible,
  keepspaces=true,
  showstringspaces=false
}

\title{\vspace{-2cm}\textbf{SkillOpt Reference}\\[0.2em]
\large Skill format, prompt templates, and implementation hyperparameters}
\author{AgentIA Engineering}
\date{\today}

\begin{document}
\maketitle
\thispagestyle{fancy}

\begin{keybox}
\textbf{Purpose.} This is the implementation companion to the proposal.
It reproduces the skill document format, the published prompt templates
from SkillOpt (arXiv:2605.23904, Appendix~C.2), the CODESKILL skill
examples (arXiv:2605.25430, Figures~4--14), and the hyperparameters
required to build the loop. Templates are verbatim from the sources so
they can be copied directly into the codebase under
\texttt{backend/app/resources/prompts/skillopt/}.
\end{keybox}

\tableofcontents
\newpage

\section{Optimization in text space}

\begin{keybox}
\textbf{The one-sentence version.} Take the four things that make weight
training reliable --- a bounded step size, a validation gate, negative
feedback, and momentum --- and apply them to a text document instead of
to a matrix of numbers. Everything else in this reference is a detail
underneath that sentence.
\end{keybox}

\subsection{Why it is not obvious}

Before SkillOpt, there were two ways to improve an LLM's guidance text.

\begin{description}[leftmargin=1.2em]
  \item[Hand-editing.] A human writes a prompt, tests it, edits it,
  tests again. Works, but does not scale, and the human's judgment is
  the only signal.
  \item[Self-editing.] The model reads its own failures and rewrites
  its own prompt. Scales, but is unreliable --- no bound on how much it
  rewrites, no validation that the rewrite helped, no memory of what
  was tried and failed. In practice, self-editing drifts: the document
  grows longer, more contradictory, and eventually worse.
\end{description}

SkillOpt is a third option: self-editing \emph{with the discipline of
an optimizer}. Bounded edits, validated against held-out data, with
negative feedback and momentum. It borrows every control that makes
gradient descent reproducible, and applies those controls to text.

\subsection{The deep-learning analogy}

The analogy is operational, not decorative. Every row matters.

\begin{table}[h]
\centering
\small
\begin{tabular}{@{}lll@{}}
\toprule
Deep learning & SkillOpt & What it prevents \\
\midrule
Weights & The skill document & (the trainable object) \\
Forward pass & Rollout batch on training tasks & Editing from too little evidence \\
Loss & Task score (0/1 from verifier) & (the objective) \\
Gradient & Reflection --- LLM proposes edits & (the proposed direction) \\
Learning rate & Edit budget \(L_t\) per step & Unbounded rewriting \\
Optimizer step & Apply bounded edits \(\to\) candidate & (the move) \\
Validation set & Held-out selection split & Overfitting to motivating examples \\
Acceptance rule & Strict improvement, ties rejected & Silent drift and degradation \\
Momentum & Epoch-wise slow/meta update & Long-horizon drift \\
Rejected gradients & Rejected-edit buffer & Repeated bad proposals \\
\bottomrule
\end{tabular}
\end{table}

The analogy is the argument: if you believe in gradient descent, you
should believe in bounded text editing with a validation gate, because
they are the same algorithm in different spaces.

\subsection{Why the gate is inside the update}

This is the part a technical judge will care about. In standard SGD:

\begin{promptbox}
weights \(\leftarrow\) weights \(-\) lr \(\cdot\) gradient
\hfill (always take the step)
\end{promptbox}

You always move. Validation runs separately to pick checkpoints.

In SkillOpt:

\begin{promptbox}
candidate \(\leftarrow\) apply\_bounded\_edits(skill)\\
if score(candidate) \(>\) score(current):\\
\hspace*{2em}skill \(\leftarrow\) candidate\\
else:\\
\hspace*{2em}reject, log to buffer
\end{promptbox}

\textbf{The step is conditional on validation.} The gradient (the
LLM's reflection) proposes; the held-out score disposes. This is closer
to a \emph{trust-region} or \emph{line-search} method than to vanilla
SGD --- and it is necessary because the ``gradient'' here is a language
model's guess, not a computed direction. A language model's reflection
can look plausible and be wrong. The gate is what makes an unreliable
proposer safe.

The consequence: the loop is monotone. The skill's selection score
never goes down, by construction. Every accepted edit is a measurable
improvement on held-out data.

\subsection{The three failure modes of naive text editing}

If someone pushes back with ``isn't this just prompt engineering,''
answer with the three ways naive editing fails --- each corresponding
to a control SkillOpt adds.

\begin{enumerate}[leftmargin=1.4em]
  \item \textbf{Unbounded editing.} Naive self-editing rewrites the
  whole document. You lose the ability to attribute improvement to any
  specific change, and the document drifts. SkillOpt bounds each step
  to \(L_t\) edits (default \(4\)), preserving continuity between
  versions.
  \item \textbf{Unvalidated editing.} Naive self-editing accepts any
  plausible rewrite. A plausible textual diagnosis can still hurt the
  actual model. SkillOpt requires strict improvement on a held-out
  split before any candidate becomes current.
  \item \textbf{No momentum.} Naive self-editing only compares adjacent
  versions. It cannot see slow drift across many steps. SkillOpt re-runs
  the \emph{same tasks} across epoch boundaries and writes durable
  guidance to a protected section that step-level edits cannot touch.
\end{enumerate}

The paper's ablation is the empirical proof: removing the momentum
mechanisms (meta skill + slow update) drops SpreadsheetBench from
\(77.5\) to \(55.0\) --- a \(22.5\)-point degradation. That is the cost
of skipping the control.

\subsection{Why ``text space'' is the right phrase}

The optimization is literally happening over the space of possible
skill documents. Each document is a point in that space. Each edit is a
move. The validation score is the objective. The gate keeps you on the
improving side.

The subtlety: this space is \textbf{non-differentiable and open-ended}.
You cannot take a gradient, and the set of possible edits is anything a
language model can write. SkillOpt replaces the gradient with a
language-model reflection, and replaces the step-size schedule with an
edit-count budget. Everything else --- validation, momentum, negative
feedback --- is directly imported.

\subsection{How to say it at different depths}

\begin{description}[leftmargin=1.2em]
  \item[Ten seconds.] ``It's a training loop for a text document.
  Each round, the model proposes small edits, and the edit is only
  kept if it measurably improves performance on a held-out set. Same
  discipline as training a neural network --- just applied to a prompt
  instead of to weights.''
  \item[Thirty seconds.] ``You have a skill document that gets
  prepended to the generation prompt. An optimizer model reads rollout
  trajectories and proposes atomic edits --- append, insert, replace,
  delete --- capped at a fixed number per step. Each candidate document
  is scored on a held-out split, and accepted only if the score
  strictly improves. Rejected edits go into a buffer so the next round
  does not repeat them. At each epoch, the same tasks are re-run under
  the old and new documents, and durable guidance gets written to a
  protected section. The controls --- learning rate, validation,
  momentum, negative feedback --- are exactly the ones you would want
  in SGD.''
  \item[Two minutes.] The four-control mapping above, plus the
  trust-region framing of the gate, plus the \(22.5\)-point ablation
  as evidence that the controls matter.
\end{description}

\begin{keybox}
\textbf{The three sentences to have ready.}
(1) \emph{The skill document is the trainable parameter, and the model
stays frozen.} No weights touched, no fine-tuning, no training
infrastructure --- just text.
(2) \emph{The step is validated before it is taken.} Every candidate
must strictly beat the current version on a held-out set. This is what
makes self-editing safe where naive self-editing drifts.
(3) \emph{The four controls --- bounded edits, validation gate,
rejected-edit buffer, epoch-wise momentum --- are exactly the controls
that make weight-space optimization reproducible, applied to text
space.}
\end{keybox}

\section{Skill document format}
\label{sec:format}

Every skill is a standalone markdown file. This format is the
convergence point of all three 2026 papers (SkillOpt, CODESKILL,
SkillBrew); it is the contract between authoring, injection, and
optimization.

\begin{promptbox}
\# Title\\
\# <short reusable name --- no repo, class, file, or variable names>\\[0.4em]
\#\# Granularity\\
task-level | event-driven\\[0.4em]
\#\# When to apply\\
<transferable trigger condition --- high-level, not instance-specific>\\[0.4em]
\#\# Rules\\
1. <actionable rule>\\
2. <actionable rule>\\
...\\[0.4em]
<!-- SLOW\_UPDATE\_START -->\\
<!-- SLOW\_UPDATE\_END -->
\end{promptbox}

\subsection{Field semantics}

\begin{itemize}[leftmargin=1.4em]
  \item \textbf{Title.} Short, reusable, capability-oriented.
  Describes \emph{what the skill enables}, not a specific instance of it.
  \item \textbf{Granularity.} Two values. \emph{Task-level} for
  whole-task strategies (e.g., ``diagnose Maven build failures'').
  \emph{Event-driven} for local trigger-response patterns
  (e.g., ``handle missing test-jar dependency'').
  \item \textbf{When to apply.} The trigger condition. High-level
  enough to fire on unseen tasks of the same family; specific enough
  not to fire on every task. This is the hardest field to author well.
  \item \textbf{Rules.} Ordered, procedural, actionable. Each rule
  should describe a concrete action or check. Prefer numbered steps
  over prose.
  \item \textbf{Protected section.} Empty in seed skills. Only the
  epoch-wise slow/meta update (see \S\ref{sec:slow}) may write here.
  Step-level edits are blocked from touching it by design.
\end{itemize}

\subsection{Authoring rules}

\begin{enumerate}[leftmargin=1.4em]
  \item Rules must be \textbf{procedural}, not declarative.
  \item Rules must be \textbf{ordered} when sequencing matters.
  \item Trigger conditions must be \textbf{transferable}.
  \item \textbf{No instance-specific identifiers.} No repo names, file
  paths, class names, module names, or one-off literals.
  \item Every rule must be \textbf{grounded} in observed failure
  evidence or constitutional principles.
  \item Combined length per skill: \textbf{under 500 tokens}.
  \item Combined length of all seed skills: \textbf{under 2,000 tokens}.
\end{enumerate}

\section{Reference skill examples}
\label{sec:examples}

Two examples from CODESKILL's maintained bank. They are Java/Maven
specific and represent the target quality bar for AgentIA's seed skills.

\subsection*{Task-level example}

\begin{promptbox}
\# Diagnose Maven build failures via environment, config, and dependency
resolution.\\[0.3em]
\#\# Granularity\\
task-level\\[0.3em]
\#\# When to apply\\
When a Java/Maven project fails to build and the cause may involve
toolchain or wrapper configuration, repository settings, missing or
stale dependencies, or build-phase ordering rather than source-code
errors.\\[0.3em]
\#\# Rules\\
1. Identify the project type and build instructions from repository
metadata before making changes.\\
2. Run the build with the project's Maven wrapper first, then capture
both early and later output to distinguish dependency-download progress
from the actual failure.\\
3. Use the first concrete failure message to distinguish network or
repository issues, missing artifacts, and compilation errors before
changing code.\\
4. If the build fails in the enforcer or toolchain phase, use verbose
output to identify the required Java and Maven versions.\\
5. Check active Java and Maven toolchain versions and align them with
the project's documented requirements.\\
6. When a dependency cannot be resolved, check whether the missing
artifact is generated later in the build.\\
7. Inspect local Maven settings and repository configuration before
changing code.\\
8. After adjusting toolchain, wrapper, repository, or dependency
configuration, rerun and re-check the logs to confirm the failure mode
changed.
\end{promptbox}

\subsection*{Event-driven example}

\begin{promptbox}
\# Handle missing Maven test-jar dependencies.\\[0.3em]
\#\# Granularity\\
event-driven\\[0.3em]
\#\# When to apply\\
When a Maven build fails with a dependency resolution error for a
test-jar or similar scoped artifact that has not yet been produced or
installed.\\[0.3em]
\#\# Rules\\
1. Inspect the failing module's POM to confirm dependency type and
scope.\\
2. Check the producing module's POM for an existing execution for the
required artifact.\\
3. If the test jar is only generated when tests run, either run the
tests or explicitly invoke the Maven test-jar goal before rebuilding.\\
4. If the artifact is missing, add or configure the appropriate plugin
execution in the producing module's POM.\\
5. Do not retry the same Maven install command with tests skipped;
skipping prevents the test jar from being created.\\
6. After generating or configuring the artifact, run a full install
so it is installed into the local repository.\\
7. Verify the artifact is present in the local Maven repository before
retrying the original command.
\end{promptbox}

\section{SkillOpt prompt templates}
\label{sec:prompts}

All templates below are reproduced verbatim from SkillOpt, Appendix~C.2,
with terminology aligned to the optimizer/target framing. They are
designed to be dropped into \texttt{prompts/skillopt/} as-is. Prompt
outputs are JSON so the pipeline can parse, filter, apply, and validate
without manual intervention.

\subsection{Failure analysis --- \texttt{analyst\_error.md}}

\begin{promptbox}
You are an expert failure-analysis agent for AI agent tasks.

You will be given MULTIPLE failed agent trajectories from a single
minibatch and the current skill document. Your job is to identify the
most important COMMON failure patterns across the batch and propose a
concise set of skill edits.\\[0.4em]
\#\# Analysis Process\\
1. Read ALL trajectories in the minibatch.\\
2. Identify the most prevalent, systematic failure patterns.\\
3. For each pattern, classify its failure type.\\
4. Propose skill edits that address the COMMON patterns, not individual
edge cases.\\
5. Edits must be generalizable; do not hardcode task-specific values.\\
6. Only patch gaps in the skill; do not duplicate existing content.\\[0.4em]
You will be told the maximum number of edits (the budget L). Produce
AT MOST L edits.\\[0.4em]
Respond ONLY with a valid JSON object:\\
\{\\
\hspace*{1em}"reasoning": "<why these edits address the batch's common
failures>",\\
\hspace*{1em}"edits": [\\
\hspace*{2em}\{"op": "append", "content": "<markdown to add at end of
skill>"\},\\
\hspace*{2em}\{"op": "insert\_after", "target": "<exact heading/text>",
"content": "<markdown>"\},\\
\hspace*{2em}\{"op": "replace", "target": "<exact text>", "content":
"<replacement>"\},\\
\hspace*{2em}\{"op": "delete", "target": "<exact text>"\}\\
\hspace*{1em}]\\
\}
\end{promptbox}

\subsection{Success analysis --- \texttt{analyst\_success.md}}

\begin{promptbox}
You are an expert success-pattern analyst for AI agents.

You will be given MULTIPLE successful agent trajectories from a single
minibatch and the current skill document. Identify generalizable
behavior patterns COMMON across the batch worth encoding in the
skill.\\[0.4em]
\#\# Rules\\
- Only propose patches for patterns NOT already covered.\\
- Focus on patterns appearing across MULTIPLE trajectories.\\
- Patterns must generalize beyond specific tasks.\\
- Prefer reinforcing existing sections over adding new ones.\\[0.4em]
Produce AT MOST L edits. Output: same JSON schema as failure analysis
with a "success\_patterns" array added.
\end{promptbox}

\subsection{Failure merge --- \texttt{merge\_failure.md}}

\begin{promptbox}
You are a skill-edit coordinator. Merge multiple independently-proposed
patches from FAILURE analysis into ONE coherent patch.\\[0.4em]
Merge guidelines:\\
1. Deduplicate: keep the best-worded version of similar edits.\\
2. Resolve conflicts: choose the one with stronger justification or
synthesize both.\\
3. Preserve unique insights: include all non-redundant corrective
edits.\\
4. Prevalent-pattern bias: edits appearing across multiple patches get
HIGH priority.\\
5. Independence: no two edits may target the same text region.\\
6. Support count: for each merged edit, estimate how many source
patches support it.\\
7. PROTECTED SECTION: do not target content between
\texttt{SLOW\_UPDATE\_START} and \texttt{SLOW\_UPDATE\_END}.
\end{promptbox}

\subsection{Success merge --- \texttt{merge\_success.md}}

Same structure as failure merge, applied to success-driven patches. Be
conservative: success patches reinforce existing behavior; only include
edits for patterns not already in the skill.

\subsection{Cross-priority merge --- \texttt{merge\_priority.md}}

\begin{promptbox}
Merge failure and success patches.\\[0.4em]
1. FAILURE PATCHES TAKE PRIORITY: failure-driven edits are preserved
unless they directly conflict with a well-supported success pattern.\\
2. Deduplicate: when a failure and success edit cover the same point,
keep the failure version.\\
3. Preserve success insights for patterns not addressed by failures.\\
4. Higher-level merges represent broader consensus; give priority.\\
5. Carry forward support\_count and source\_type for each edit.\\
6. PROTECTED SECTION: do not target protected content.
\end{promptbox}

\subsection{Ranking and selection --- \texttt{ranking.md}}

\begin{promptbox}
You are an expert edit-ranking optimizer for a skill optimization
system. Rank the edits and select the top ones.\\[0.4em]
Ranking criteria (in priority order):\\
1. Systematic impact: edits addressing widespread, recurring failure
patterns rank highest. A rule fixing 50\% of failures beats one fixing
a single edge case.\\
2. Complementarity: edits filling gaps, not duplicating content.\\
3. Generality: general principles rank higher than instance-specific
advice.\\
4. Actionability: concrete guidance ranks higher than vague advice.\\[0.4em]
Respond ONLY with:\\
\{"reasoning": "...", "selected\_indices": [<0-based indices>]\}
\end{promptbox}

\subsection{Epoch-wise slow update --- \texttt{slow\_update.md}}
\label{sec:slow}

\begin{promptbox}
You are a strategic skill advisor for an AI agent optimization system.
You see how the skill has evolved across an entire epoch by comparing
the SAME tasks under two consecutive skill versions. This longitudinal
view lets you identify systemic drift, regressions, and persistent
blind spots that step-level edits cannot catch.\\[0.4em]
\#\# What you receive\\
1. Previous epoch's skill and current epoch's skill.\\
2. Longitudinal comparison: same 20 training tasks under both skills,
categorized into regressions, persistent failures, improvements, and
stable successes.\\
3. Previous slow update guidance, if any.\\[0.4em]
\#\# Your process\\
1. Reflect on the previous guidance: what was effective, what
backfired, what was missed.\\
2. Write updated guidance that retains what worked, revises what
didn't, adds new instructions for observed regressions and persistent
failures.\\[0.4em]
\#\# Output requirements\\
Write a strategic guidance block that will OVERWRITE the previous
guidance in the protected section. Prioritize: (1) preventing
regressions, (2) fixing persistent failures, (3) reinforcing successful
patterns. Do NOT duplicate content already in the main skill body.\\[0.4em]
Respond ONLY with:\\
\{"reasoning": "...", "slow\_update\_content": "<exact text>"\}
\end{promptbox}

\subsection{Optimizer meta skill --- \texttt{meta\_skill.md}}

\begin{promptbox}
You are an optimizer coach. Your job is to write a compact
optimizer-side meta skill that helps future optimizer calls produce
better skill edits in this environment.\\[0.4em]
\#\# What you receive\\
1. Previous epoch's last-step skill.\\
2. Current epoch's last-step skill.\\
3. Longitudinal comparison on the SAME sampled tasks.\\
4. Previous optimizer memory, if one existed.\\[0.4em]
\#\# Your goal\\
Capture: which kinds of edits help in this environment; which tend to
be vague, redundant, brittle, or harmful; what level of abstraction
works best; which failure-repair patterns to prioritize; which
regression risks to guard against.\\[0.4em]
\#\# Constraints\\
- Address the FUTURE OPTIMIZER directly, not the training model.\\
- Use evidence from the adjacent-epoch comparison, not generic advice.\\
- Keep it compact and high-signal.\\[0.4em]
Respond ONLY with:\\
\{"reasoning": "...", "meta\_skill\_content": "<compact guidance>"\}
\end{promptbox}

\section{CODESKILL supplementary templates}
\label{sec:codeskill}

CODESKILL's prompt set (Figures~6--14) is complementary to SkillOpt's.
Use CODESKILL's extraction and maintenance prompts for the initial
skill authoring pass, and SkillOpt's prompts for the optimization loop.

\subsection{Task-level extraction --- Figure 6}

\begin{promptbox}
1. A bootstrap skill is a task-level reusable pattern supported by
multiple trajectories. It must be broader than a single local event.\\
2. Generate only if the pattern is reusable across similar tasks. Keep
\texttt{when\_to\_apply} high-level and write transferable, actionable
rules.\\
3. Ground every part of the skill in repeated evidence from the
trajectories and outcomes. Do not invent unsupported steps, checks, or
guidance. If failed trajectories reveal reusable cautions, include them
as cautionary rules.\\
4. Skip if the evidence is weak, contradictory, accidental, too local,
or collapses into an event-level reaction instead of a task-level
pattern.\\
5. Do not include repository names, issue descriptions, exact task
goals, variable names, function names, class names, module names,
exact file paths, or one-off literals.\\[0.4em]
Output schema:\\
\{"action": "generate", "skill": \{"title": "...", "granularity":
"general", "when\_to\_apply": "...", "rules": [...]\}\}\\
or \{"action": "skip", "reason": "..."\}
\end{promptbox}

\subsection{Event-driven extraction --- Figure 7}

\begin{promptbox}
You are an expert at extracting reusable memory from bash-agent
trajectories. Extract one reusable event-driven skill from a full
trajectory. Use only the provided evidence.\\[0.4em]
\#\# Rules\\
1. An event-driven skill is a local trigger-response pattern. Focus on
one important event; stay narrower than a whole-task workflow.\\
2. Generate only if the event yields a reusable lesson beyond this
exact task.\\
3. Ground the trigger and guidance in the trajectory and result. Do
not invent an event.\\
4. If multiple candidate events exist, choose the single most
reusable one.\\
5. Skip if there is no strong reusable local event, or if the lesson
is too task-specific or workflow-level.\\
6. No repository names, file paths, or one-off literals.\\[0.4em]
Output: same schema with \texttt{"granularity": "event-driven"}.
\end{promptbox}

\subsection{Evolution --- Figure 8}

\begin{promptbox}
Choose exactly one action: \texttt{evolve} (revise exactly one existing
skill) or \texttt{skip} (make no revision).\\[0.4em]
\#\# Rules\\
1. Alignment and need: choose evolve only when one existing skill is
clearly aligned with the current trajectory and the result shows it
should be revised.\\
2. Reusability: revise only if the trajectory reveals a reusable
missing case, missing check, better decision rule, improved ordering,
or reusable caution.\\
3. Applicability: \texttt{when\_to\_apply} remains high-transferable.\\
4. Grounding: the revision must be directly supported by the provided
trajectory and result.\\
5. Identity preservation: keep the same capability identity.\\
6. Skip if the evidence is weak, contradictory, or the trajectory is
better described as a brand-new skill.
\end{promptbox}

\subsection{Bank maintenance --- Figure 9}

\begin{promptbox}
Choose exactly one action: \texttt{add}, \texttt{merge}, or
\texttt{drop}.\\[0.4em]
1. \textbf{add}: candidate is reusable, coherent, and distinct from
retrieved skills.\\
2. \textbf{merge}: candidate and one retrieved skill share the same
capability identity and can be combined into one stronger skill.\\
3. \textbf{drop}: candidate is covered by stronger skills, redundant,
weakly evidenced, unsafe, too local, or too task-specific.\\
4. No resulting skill may include repo names, file paths, or one-off
literals.
\end{promptbox}

\section{Hyperparameters}

Default configuration for the SkillOpt loop. Values from Table 4 of the
paper, with AgentIA-specific adaptations noted.

\begin{table}[h]
\centering
\begin{tabular}{@{}lll@{}}
\toprule
Parameter & Value & Notes \\
\midrule
Edit budget \(L_t\) & 4 & cosine decay, floor 2 \\
Rollout batch size & 40 & scale to session volume \\
Reflection minibatch & 8 & fixed \\
Epochs & 4 & fixed \\
Slow-update samples/epoch & 20 & fixed \\
Optimizer model & frontier & different provider than target \\
Gate criterion & strict \(>\) & ties rejected \\
\bottomrule
\end{tabular}
\end{table}

\subsection{AgentIA adaptations}

\begin{itemize}[leftmargin=1.4em]
  \item \textbf{Verifier.} \texttt{mvn test -o} exit code \(+\) constitutional
  validators. The paper assumes a scalar verifier; the composite score
  is AgentIA-specific.
  \item \textbf{Selection split.} 20\% of recent sessions, never used
  for training or authoring. Minimum 10 sessions; prefer 15.
  \item \textbf{Reflection batch.} Scale the paper's 40 down to session
  volume (30 is realistic for the pilot).
  \item \textbf{Optimizer model.} A different provider via
  \texttt{LLMFactory} to avoid correlated generate-and-critique failure.
\end{itemize}

\section{Implementation checklist}

\begin{enumerate}[leftmargin=1.4em]
  \item Copy all prompt templates from \S\ref{sec:prompts} into
  \texttt{backend/app/resources/prompts/skillopt/} as separate
  \texttt{.md} files.
  \item Author seed skills using the format in \S\ref{sec:format} and the
  examples in \S\ref{sec:examples}. Save to\newline
  \texttt{backend/app/resources/skills/}.
  \item Implement the edit applier with the four atomic operations
  (append, insert\_after, replace, delete). Enforce the protected
  section guard.
  \item Implement the hierarchical merge in the order: failure merge,
  success merge, cross-priority merge.
  \item Implement the ranker that clips to \(L_t\).
  \item Implement the gate: score the candidate on the selection split,
  accept only if strictly greater than current.
  \item Persist rejected edits with score drops to the buffer.
  \item Implement the epoch-boundary slow update; write to the
  protected section.
  \item Implement the meta-skill capture; keep optimizer-side only.
  \item Wire the loop: forward pass on the training split, reflection
  minibatches, merge, rank, apply, gate, repeat.
\end{enumerate}

\begin{keybox}
\textbf{The one thing not to skip.} The epoch-wise slow/meta update
(\S\ref{sec:slow}). SkillOpt's ablation shows removing both the meta
skill and the slow update costs \(22.5\) points on SpreadsheetBench ---
the largest single degradation in the paper. Per-step reflection cannot
substitute for it.
\end{keybox}

\section{Sources}

\begin{itemize}[leftmargin=1.4em]
  \item Li et al. \emph{SkillOpt: A controllable text-space optimizer
  for agent skills.} arXiv:2605.23904 (2026). Appendix~C.2 is the
  source for all prompts in \S\ref{sec:prompts}; Table~4 for
  hyperparameters.
  \item Li et al. \emph{CODESKILL: Learning self-evolving skills for
  coding agents.} arXiv:2605.25430 (2026). Figures~4--14 for skill
  format, examples, and extraction/evolution/maintenance prompts.
  \item Hu et al. \emph{SkillBrew: Multi-objective curation of skill
  banks.} arXiv:2605.29440 (2026). Referenced for deferred bank-level
  curation.
\end{itemize}

\end{document}